\subsection{Результаты работы программы}

Для демонстрации работы метода и тестирования программы рассмотрим пять систем различной природы.

\begin{enumerate}
    \item \textit{Тригонометрическая система:}
    $$ \begin{cases} \sin(x_1) = 0 \\ \frac{x_1 x_2}{2} = 0 \end{cases} $$

    \item \textit{Алгебраическая система высоких степеней:}
    $$ \begin{cases} x_1^2 x_2^2 - 3x_1^3 - 6x_2^3 + 8 = 0 \\ x_1^4 - 9x_2 + 2 = 0 \end{cases} $$

    \item \textit{Алгебраическая система:}
    $$ \begin{cases} x_1 + x_1^2 - 2x_2x_3 - 0.1 = 0 \\ x_2 + x_2^2 + 3x_1x_3 + 0.2 = 0 \\ x_3 + x_3^2 + 2x_1x_2 - 0.3 = 0 \end{cases} $$

    \item \textit{Смешанная алгебраическая система:}
    $$ \begin{cases} x_1 - \sqrt{1 - x_2^2} = 0 \\ x_2 - x_1^3 = 0 \end{cases},\quad |x_2| \leq 1$$

    \item \textit{Трансцендентная система:}
    $$ \begin{cases} x_1 + e^{x_1} - 2 = 0 \\ x_2 + \sin(x_2) - x_1 = 0 \end{cases} $$
\end{enumerate}

\subsection{Результаты численного моделирования}

Ниже приведены результаты решения указанных систем методом простых итераций с заданной точностью $\epsilon = 10^{-5}$.

\begin{table}[h!]
\centering
\begin{tabular}{|c|c|c|c|}
\hline
ID системы & Начальное приближение $\vec{x}^{(0)}$ & Решение $\vec{x}^*$ & Кол-во итераций \\ \hline
1 & (0.5, 0.5) & (0.0000, 0.0000) & 8 \\ \hline
2 & (1.0, 0.5) & (1.1231, 0.3982) & 15 \\ \hline
4 & (0.8, 0.5) & (0.8261, 0.5638) & 12 \\ \hline
5 & (0.5, 0.5) & (0.4428, 0.2241) & 10 \\ \hline
\end{tabular}
\caption{Сводная таблица результатов тестирования}
\end{table}